
\section*{修正表示原稿の参照表について}
上の英語版から自動同期される表は、各意見に対応する改訂後の節・命題と、実際の修正表示PDFで検証したページ・余白行番号を示す。節への参照は複数の改訂段落を含みうるが、定理類の項目名は対象となる記述を特定する。

\section*{編集者への回答}
\Comment{E.1}
「査読者はともに本論文の科学的貢献に価値があると認めている。しかし両者とも提示方法に多くの問題を指摘しており、その多くは単なる文章の修正にとどまらず、論文の科学的部分にも及んでいる。学習に関する部分では特にそうである。……とくに正しさに関する問題は慎重に検討すべきである。」

\Reply
同意する。投稿時の学習器を保守的な Gold 学習器に置き換え、終端規則に関する欠落を証拠集合と完全性証明の双方で補修し、使用されていない外側文脈の注釈を削除した。固定 $h$ の代入可能性は非空の内部因子について定義し、$\lambda$ は別途扱う。第4節でバッチ作用素と逐次学習器を定義した後、第5節で完全性のための標的側の生成語型付けを導入する。多項式特徴データは集合依存バッチ作用素に関する主張であり、逐次学習器については多項式時間更新を伴う Gold の安定化を示す。
\Location{定義2.1、第4.2節、第5.1節と定理5.3、系5.5、第6--8節。}

\section*{Reviewer 1}
\subsection*{総評と学習アルゴリズム}

\Comment{R1.1}
「本投稿の主結果は、Clark and Eyraud (2007) の代入可能文脈自由言語（SCFL）の学習、および Yoshinaka (2008) の $(k,\ell)$-SCFL の学習を洗練された形で一般化したものである。……唯一の重大な誤りは学習アルゴリズムの定義にあるが、先行研究に従えば容易に修正できる。」

\Reply
中心的な正しさの問題を指摘されたことに感謝する。正例ごとにバッチ文法を再構成する旧方式では構文的安定化が保証されないため，集合依存の再構成作用素 $\Bh$（第4.1節）と保守的な逐次学習器 $\Ah$（第4.2節）を分離した。後者は現在の仮説が生成しない正例を受け取った場合にのみ再構成し，系5.5が最終的な構文的安定化を証明する。Eyraud--Heinz--Yoshinaka (2016, Appendix) の遅延再構成の現象も引用した。したがって逐次学習器に主張するのは Gold 収束と多項式時間更新であり，逐次的な多項式時間・データ同定（同 Definition 2.5）は主張しない。多項式特徴データの保証は集合依存の $\Bh$ に対するもので（同 Definition 2.8），明示した定量的範囲に限る。第2.1節でこの区別を説明した。
\Location{第2.1節、第4.2節、系5.5、定理7.2、定理8.3およびその後の補足。}

\Comment{R1.2}
「一方、提示にはかなり改善の余地がある。技術用語の定義の欠如、不必要に長い証明、メタ変数の使い方の不統一、節を過度に細分化している点を直すことで、読者の負担を大幅に減らせるはずであり、そのような大改訂を勧める。……この内容は現在の半分以下の長さで提示できるのではないか。」

\Reply
提示を大幅に短縮した。投稿時原稿は40ページであり，付録・参考文献を含む改訂稿はClean版・修正表示版ともに27ページである。説明の反復を除き，証明義務を残して序論・結論・線形正規化を短縮した。番号付き定理類は半数以下，本文の節数は11から10，小節数は31から16となった。第2節に定義と正例学習モデルをまとめ，旧計算量節を第5節に統合した。第6--8節では定量的議論を整理し，旧三つの境界関連の節は短縮した第9節に統合した。

未定義の用語とメタ変数の不統一を修正し（R1.8，R1.11），SSBNFの説明は本文に置いた。両査読者の提案に従い，型付き生成語と正規化の証明を付録A，線形spineの正規化を付録B，命題6.3(i)の証明を付録Cに配置した。
\Location{第1--2節、第4--9節、付録A--C。}

\subsection*{投稿時第1節}
\Comment{R1.3}
「研究分野内での位置づけと対象言語族の重要性を述べれば、論文の魅力が増す。一方、現状は技術的な内容に重きを置きすぎているので、より簡潔にすべきである。」

\Reply
同意する。第1.2節では固定 $h$ の定性的な学習構成と Yoshinaka の一般的枠組みとの関係を簡潔に述べ，貢献を定量的保証と構造・表現力上の位置づけの二つに整理した。固定準同型は有限の型付けバイアスであり，学習を主張する固定 $h$ のクラスと，表現力の比較に使う $\RS$ を区別した。2014年のプレプリントと Yoshinaka (2015) との関係，固定有限モノイド型付けに固有の定量的・代数的結果も明示した（R2.1を参照）。
\Location{第1節。}

\Comment{R1.4}
「$\mathcal D_L$ と $\mathcal C_h^{\mathrm{cf}}$ は標準的な記号ではないので定義が必要である。とくに『一般の文脈自由クラス $\mathcal C_h^{\mathrm{cf}}$』という言い方は誤解を招きやすい。」

\Reply
同意する。$\mathcal D_L(x)$ は序論での初出時と予備事項で定義した。$\mathcal C_h^{\mathrm{cf}}$ も序論の初出時に $\sim_h$-代入可能な文脈自由言語のクラスとして定義する。$h$ を固定した制約のないクラス $\RSh$ は、合併 $\RS=\bigcup_h\RSh$ より先に導入した。第3節でこれらのクラスと $\Chlin$ を形式的に扱い、「一般の文脈自由クラス」という表現は削除した。
\Location{序論、第2.2節および第3節冒頭。}

\Comment{R1.5}
「貢献が(1)--(6)と列挙されているが、細かな技術的事項まで独立の貢献として数える必要はない。また、この段階では専門用語の意味が未だ不明なものが多く、読者はついていけない可能性がある。」

\Reply
同意する。第1.2節では定性的学習構成と Yoshinaka の枠組みを番号なしの短い段落で紹介し，番号付き貢献を定量的保証と構造・表現力上の位置づけの二つに絞った。定量的貢献では特徴データの多項式上界と指数下界を概説し，「型付き厚み」と「開始記号分離二分標準形」に言及するが，それぞれの形式的定義は第6節と第5.1節で与える。構造上の貢献では固定窓の判定条件を古典的結果として示し，「局所自明」は第2.2節で定義した。第1.1節の $L=\{a,aa\}$ の例は，自明型付けでは代入可能性が破れる一方，偶奇型付けで因子を区別できることを示す。これは定義の例であり，固定窓クラスからの分離例ではない。
\Location{第1.1節（固定 $h$ 代入可能性の定義直後の最小例）、第1.2節。}

\Comment{R1.6}
「『Yoshinaka の古典的言語』は、文献では『Łukasiewicz 言語』と呼ぶことが多い。」

\Reply
同意する。「Łukasiewicz 言語」を用い、Autebert--Berstel--Boasson (1997) のハンドブック文献を加え、1種類の括弧からなる Dyck 言語との関係を簡潔に述べた。この言語がすべての固定 $(k,\ell)$ クラスの外にあるという既知の観察は Yoshinaka (2008, Section 3) に帰属させ、有限型付けの合併 $\RS$ の外にあるという本論文の強い結論と区別した。
\Location{第9.3節。}

\subsection*{投稿時第2節}
\Comment{R1.7}
「『この論文で考える学習器は、正例の提示順序によらず同じ出力を生じる』という主張は疑わしい。学習アルゴリズムと、その部分手続きである文法構成アルゴリズムが混同されているように見える。実際、第5節の学習アルゴリズムは収束せず、Gold の基準を満たさない。Clark と Eyraud に従って修正しても、例えば $a^ncb^n$ の学習では、特徴標本を最初に受け取る場合と、最初の正例が非常に長い場合とで、最終出力は異なるだろう。」

\Reply
同意する。投稿時の順序独立性の主張を撤回し，集合依存のバッチ作用素と保守的逐次学習器を区別した。第4.2節では任意の無限標的について，十分長い正例をテキストの先頭に追加すると，安定後にも余分な観測非終端記号が残りうることを示し，提示順序によって最終文法が異なると説明する。Clark (2013, Section 4, Example 1) にも同様の例がある。
\Location{第2.1節、第4.2節、系5.5。}

\Comment{R1.8}
「この節では、（線形）文脈自由文法を含め、必要な記法と用語を十分に定義すべきである。『type』『reconstruction』『realize』『reachable』『observation』『control』『local configuration』などが脚注で場当たり的に定義されている。数学的定義のない直感的な言葉に過ぎないなら、名前を付けないほうがよい場合もある。過剰な用語の氾濫は読者の負担になり、理解を妨げる。」

\Reply
同意する。第2節に線形CFG，到達可能・生成可能・縮約済み文法，正例学習モデル，再構成，特徴標本，分布，固定 $h$ のクラスをまとめた。証明内部の不要な術語を削除し，部分語に対応する仮説非終端記号 $\NT{x}$ は第4.1節で初出時に定義した。文字数 $|w|_a$，SGL と $(k,\ell)$-SGL の略称，$\mathsf{Acc}\subseteq M$ も初出時に説明し，未使用の $[x]_{\equiv_L}$ を削除した。列挙された語のうち，\emph{reachable} と生成語の \emph{type} は各定義箇所で，\emph{reconstruction} はバッチ作用素として定義する。\emph{observation} は独立の術語にせず，\emph{realize} と \emph{local configuration} は削除した。\emph{control} は制御集合などの標準的な用法に限る。
\Location{第2節、第4.1節、第5.1節、付録A.2。}

\Comment{R1.9}
「『finite typed reconstruction basis』という語はここでは定義されているが、その四つの単語が個別には未定義なのに全体が不可分の概念なら、もっと短い名称のほうがよい。」

\Reply
同意する。名称を変えるのではなく、証明内部のこの対象と『基底』自体を削除した。新しい完全性証明は、生成語型付きの標的文法と厳密な再構成に必要な有限証拠集合 $\mathsf{Wit}(\widetilde G)$ のみを用いる。
\Location{第5.1節。}

\subsection*{投稿時第3節}
\Comment{R1.10}
「論文は $h$ が固定されることを、やや過度かつ不自然な形で強調している。それに合わせるなら $\mathsf{RS}_h$ を $\mathsf{RS}$ より先に定義すべきだ。$\mathsf{RS}$ が正例から学習可能だとは示されておらず、実際それは学習可能でない。学習可能なのは $\mathsf{RS}_h$ である。」

\Reply
同意する。序論で $\RSh$ を $\RS:=\bigcup_h\RSh$ より先に定義し、Gold 同定定理が扱うのは、固定 $h$ に対する $\Chcf=\RSh\cap\mathrm{CFL}$ の非空言語であって、$\RSh$ 全体でも変動する $h$ に関する合併 $\RS$ でもないことを明示した。後者は表現力の比較にのみ使う。また、$\RS$ と $\RS\cap\mathrm{CFL}$ は非空アルファベット上で正例から同定可能ではない。すべての有限言語と $\Sigma^*$ を含むためである。
\Location{要旨、序論、第2.2節および第3節冒頭。}

\Comment{R1.11}
「定義3.2は注意を要する。Clark and Eyraud (2007) の定義の述べ方はいくらか不注意だが、同論文第6.1節で $a^+$ が代入可能であると述べ、文法構成アルゴリズムで空語を除外していることを考えると、Yoshinaka (2008) の $(0,0)$-代入可能性を意図したと解するのが妥当だろう。……」

\Reply
査読者による Clark--Eyraud (2007) の解釈に同意する。非空内部因子に制限すれば意図された場合は Yoshinaka の $(0,0)$ と一致する（第3.1節）。Clark (2013) を引用し，Clark--Eyraud の Definition 5 が一般の文字列を量化する点と，改訂稿の非空因子の規約を区別した。固定 $h$ の代入可能性は非空の内部因子に課し，$\lambda$ は開始記号で扱う。例えば $L=a^+$ では $(a,\lambda)$ は $\mathcal D_L(\lambda)$ と $\mathcal D_L(a)$ の双方に属するが，$(\lambda,\lambda)$ は後者だけに属するので，空因子の比較は $a^+$ を排除する。中心印付き言語の例でも空因子の比較は用いない。$(k,\ell)$ の定義では外側の文脈を $u_1,u_2,v_1,v_2$，固定窓を $p,q$，可変内部を $r,s$ とした。
\Location{定義2.1、第3.1節の Clark--Eyraud の特殊例および固定窓の定義。}

\Comment{R1.12}
「命題3.5：$\sqcup$ と $\mathrm{pref}_k$ が定義されていない。……証明が非常に長い。$T$ の積を最初に明示すれば、$\tau$ がモノイド準同型であることは一行で済む。$h_{k,\ell}$ や $M_{k,\ell}$ を使った言い換えは不必要に煩雑である。」

\Reply
固定窓の議論が冗長だったことに同意する。未定義の $\sqcup$ を削除し，接頭辞・接尾辞と短語の規約を先に定義した。命題3.3は有限像の積を直接与え，準同型性を簡潔に確認する。ただし $h_{k,\ell}$ と $M_{k,\ell}$ は後続の命題3.5，補題7.1，定理7.2，補題7.4で再利用するため保持した。短語は個別の型を持つ。固定窓の反例判定条件は命題3.3の後にまとめ，古典的な局所自明性の特徴づけ，窓・群の分解，厚みに基づく上界は命題3.5，補題7.4，系7.5として区別した。
\Location{命題3.3と直後の固定窓反例条件、命題3.5、補題7.4、系7.5。}

\Comment{R1.13}
「$h$ がデータとして与えられるというより、早い段階で『効率的に計算可能な $h:\Sigma^*\to M$ を固定する』などと述べるほうが明確だ。（線形）文脈自由言語を定義してほしい。……『明示的有限モノイド準同型』は定義されていない。……定理の主張では技術的に不要な言葉を削除すべきである。」

\Reply
同意する。有限モノイドは積表と文字像で固定し、$h(w)$ の計算方法を明示した。線形 CFG を予備事項で定義し、主定理では「固定された効率的に計算可能な準同型」という表現を用いる。
\Location{第2--3.2節。}

\Comment{R1.14}
「『The proof for fixed $h$』という表現が理解できない。『定理3.7の証明』という意味か。」

\Reply
同意する。この紛らわしい句は削除し，改訂稿では対応する定理と証明を直接記した。
\Location{第3.2節。}
